% ============================================================
% PUNTO 3: Filtrado espacial de imágenes con núcleos 2D
% ============================================================

\subsection{Punto 3: Filtrado espacial de una imagen con núcleos 2D}

\subsubsection{Descripción del código}

El código adjunto (\texttt{cod01.py}) carga una imagen, la convierte a
escala de grises y le aplica convolución discreta 2D mediante
\texttt{cv2.filter2D} con núcleos fijos de $5\times5$: un núcleo
promediador (\texttt{kernel\_pasa\_bajas}, todos los coeficientes iguales
a $0.04$), dos núcleos con centro positivo fuerte y entorno negativo
(\texttt{kernel\_pasa\_altas00} y \texttt{kernel\_pasa\_altas02}) y una
cascada pasa-bajas seguida de pasa-altas denominada ``Filtro Pasa
Bandas''. Como el código original usa \texttt{cv2.imshow}, se adaptó a un
entorno sin ventana gráfica guardando cada resultado como imagen, sin
modificar los núcleos ni las llamadas a \texttt{filter2D}. Se empleó como
entrada una fotografía real de una tarjeta débito sobre una mesa de
madera (Fig.~\ref{fig:t2p3_orig}), que aporta bordes nítidos (logotipo,
chip, contorno de la tarjeta), regiones planas (el fondo amarillo) y
textura fina (la veta de la madera).

Las sumas de coeficientes, calculadas numéricamente y verificadas con
pruebas automáticas, son: \texttt{kernel\_pasa\_bajas} $=1.0$,
\texttt{kernel\_pasa\_altas00} $=1.0$, \texttt{kernel\_pasa\_altas02}
$=0.0$ y el núcleo nuevo del inciso (b) $=0.0$. Nótese que
\texttt{kernel\_pasa\_altas00}, pese a su nombre, no suma cero.

\begin{figure}[!t]
\centering
\begin{minipage}{0.48\columnwidth}
\centering
\includegraphics[width=\linewidth]{punto3_00_original.png}
\end{minipage}\hfill
\begin{minipage}{0.48\columnwidth}
\centering
\includegraphics[width=\linewidth]{punto3_01_pasa_bajas.png}
\end{minipage}
\caption{Imagen de entrada en escala de grises (izquierda) y resultado del filtro pasa-bajas de $5\times5$ (derecha): la veta de la madera y los bordes se suavizan, pero el brillo medio se conserva.}
\label{fig:t2p3_orig}
\end{figure}

\begin{figure}[!t]
\centering
\begin{minipage}{0.48\columnwidth}
\centering
\includegraphics[width=\linewidth]{punto3_02_pasa_altas00.png}
\end{minipage}\hfill
\begin{minipage}{0.48\columnwidth}
\centering
\includegraphics[width=\linewidth]{punto3_03_pasa_altas02.png}
\end{minipage}
\caption{Resultado de \texttt{kernel\_pasa\_altas00} (izquierda, suma $=1$) y \texttt{kernel\_pasa\_altas02} (derecha, suma $=0$).}
\label{fig:t2p3_altas}
\end{figure}

\begin{figure}[!t]
\centering
\begin{minipage}{0.48\columnwidth}
\centering
\includegraphics[width=\linewidth]{punto3_04_pasa_bandas.png}
\end{minipage}\hfill
\begin{minipage}{0.48\columnwidth}
\centering
\includegraphics[width=\linewidth]{punto3_05_kernel_nuevo.png}
\end{minipage}
\caption{Cascada pasa-bajas más \texttt{kernel\_pasa\_altas02} (izquierda) y núcleo nuevo de suma cero, Laplaciano de $3\times3$ (derecha).}
\label{fig:t2p3_nuevo}
\end{figure}

\subsubsection{(a) ¿Se puede considerar un filtro?}

Sí. La convolución con un núcleo fijo es, por definición, un filtro lineal
e invariante al desplazamiento: la operación
$y[m,n]=\sum_i\sum_j h[i,j]\,x[m-i,n-j]$ es la misma convolución de un
filtro FIR unidimensional, aplicada sobre dos ejes. Cada núcleo es una
respuesta al impulso de soporte finito, equivalente al vector de
coeficientes de un FIR. El núcleo pasa-bajas (suma $1$, todos positivos)
es un promedio móvil bidimensional: atenúa las frecuencias espaciales
altas (bordes y textura fina) y conserva el brillo local, como se observa
en la Fig.~\ref{fig:t2p3_orig}. Los núcleos con centro positivo y entorno
negativo se comportan como un Laplaciano discreto: responden donde el
brillo cambia rápido (contorno de la tarjeta, letras, chip) y poco en
regiones uniformes. \texttt{kernel\_pasa\_altas02} (Fig.~\ref{fig:t2p3_altas},
derecha) deja el fondo amarillo casi negro y solo resalta los bordes,
mientras que \texttt{kernel\_pasa\_altas00} (izquierda) conserva la imagen
completa y añade realce de detalle, porque su suma es $1$ y por tanto aún
deja pasar la componente de continua. La cascada (Fig.~\ref{fig:t2p3_nuevo},
izquierda) actúa como un pasa-bandas: primero elimina el ruido de muy alta
frecuencia y luego extrae los bordes, produciendo líneas más limpias y
tenues que \texttt{kernel\_pasa\_altas02} solo.

\subsubsection{(b) Cambio de la matriz con suma cero}

Se definió un núcleo nuevo, cuadrado y de suma cero, el Laplaciano
discreto de $3\times3$:
\begin{equation}
h=\begin{bmatrix}1&1&1\\1&-8&1\\1&1&1\end{bmatrix},\qquad
\sum_{i,j}h[i,j]=0.
\end{equation}
Un núcleo de suma cero no tiene ganancia en continua: sobre una región
uniforme de $128$ niveles de gris la salida de \texttt{filter2D} es
exactamente $0$ en todo el interior (verificado con una prueba
automática), sin importar el brillo original, y sobre un borde abrupto la
respuesta es grande ($|y|>20$). El núcleo se convierte así en un detector
de bordes puro que descarta la información de brillo. Comparado sobre la
misma región plana, \texttt{kernel\_pasa\_altas00} (suma $1$) devuelve
$128$, mientras que \texttt{kernel\_pasa\_altas02} y el núcleo nuevo
(suma $0$) devuelven $0$. En la imagen real (Fig.~\ref{fig:t2p3_nuevo},
derecha) el fondo y el interior de la tarjeta quedan negros y solo
permanecen los contornos, más finos y tenues que con el núcleo de
$5\times5$ porque el soporte es menor. Al cambiar los valores del núcleo
manteniendo suma cero, las regiones planas siempre desaparecen; lo que
cambia es la intensidad y el patrón con que responden los bordes. Además,
como \texttt{filter2D} trabaja con salida \texttt{uint8}, las respuestas
negativas se saturan a $0$ y solo se visualiza la parte positiva de la
respuesta al borde.
